\chapter{Futuras mejoras} \label{cap:futuras_mejoras}

\section{Introducción}
Aunque GuardiApp es una solución funcional y completa, el entorno hospitalario es dinámico y siempre existen oportunidades para evolucionar el sistema. En este capítulo se proponen varias líneas de mejora para futuras versiones.

\section{Propuestas de Evolución}

\begin{itemize}
    \item \textbf{Aplicación Móvil Nativa}: Desarrollar una versión nativa (iOS/Android) utilizando React Native para mejorar la experiencia de usuario en los fichajes y las notificaciones push.
    \item \textbf{Inteligencia Artificial para Cuadrantes}: Implementar un algoritmo de aprendizaje automático que sugiera los cuadrantes mensuales optimizados, teniendo en cuenta las preferencias de los médicos, descansos obligatorios y equidad en el reparto de festivos.
    \item \textbf{Integración con Calendarios Externos}: Ampliar la sincronización actual de Google Calendar a otros servicios como Microsoft Outlook o Apple Calendar.
    \item \textbf{Módulo de Estadísticas Avanzadas}: Crear un panel de control para la dirección del hospital con métricas sobre horas trabajadas, absentismo y costes de guardias por especialidad.
    \item \textbf{Sistema de Firma Digital}: Integrar la firma electrónica oficial para la aprobación de los informes PDF generados diariamente.
\end{itemize}

\section{Escalabilidad}
El sistema está preparado para escalar a nivel de infraestructura (uso de contenedores Docker o servicios Cloud como AWS/Azure), lo que permitiría dar soporte a múltiples centros hospitalarios de forma simultánea (arquitectura multi-tenancy).
